\documentclass[10pt,a4paper]{amsart}
\usepackage[margin=19mm]{geometry}
\usepackage{amsmath,amssymb,mathtools}
\usepackage[scr=rsfs,cal=euler]{mathalfa}
\usepackage{xfrac}
\usepackage[unicode,hidelinks,bookmarks=false]{hyperref}
\newcommand{\ie}{i.e.\ }
\newcommand{\cf}{cf.\ }
\newcommand{\resp}{respectively\ }
\newcommand{\Subsection}{\subsection}
\providecommand{\nobreakdash}{\nobreak\hbox{-}}
\setlength{\emergencystretch}{2em}
\multlinegap0pt
\usepackage{bm}
\usepackage{bbm}
\usepackage{dsfont}
\usepackage{xspace}
\usepackage{cancel}
\usepackage{mathtools}
\usepackage{amsthm}
\makeatletter
\@ifundefined{lemma}{\newtheorem{lemma}{Lemma}}{}
\makeatother
\newcommand{\E}{\mathcal{E}}
\newcommand{\NN}{\mathbb{N}}
\newcommand{\pl}{\hspace{.1cm}}
\renewcommand{\L}{\mathcal{L}}
\title[Self-restricting Noise and Exponential Relative Entropy Deca]{Self-restricting Noise and Exponential Relative Entropy Decay Under Unital Quantum Markov Semigroups}
\author{Nicholas LaRacuente}
\date{Source: arXiv:2203.03745v5 (CC BY 4.0)}
\begin{document}
\maketitle

\noindent\emph{Source and licence.} Self-restricting Noise and Exponential Relative Entropy Decay Under Unital Quantum Markov Semigroups. Version arXiv:2203.03745v5 (CC BY 4.0), distributed under the Creative Commons Attribution 4.0 International licence, as recorded in the arXiv licence metadata for this exact version and shown on the abstract page https://arxiv.org/abs/2203.03745v5. Full rights notice: licenses/arxiv-2203.03745-CC-BY-4.0.txt.\par\smallskip

\noindent\textit{Editorial scope.} Counted unit: the Lemma whose source label is \texttt{lem:zenolem} in the article (source0.tex lines 1234--1244, source label \texttt{lem:zenolem}), delivered together with the article's own complete original proof of it (source0.tex lines 1245--1320, one \texttt{proof} environment of 76 lines). No mathematics is written, repaired or re-derived: statement and proof are the article's own text line for line. Required context retained uncounted: none. Every definition, display and label of those lines is retained unchanged. Omitted from this package and not redistributed: 1--1233, 1321--1413. Nothing inside a retained excerpt is omitted or altered: every retained body is delivered exactly as the article's lines are written, so no omission receipt of any kind accompanies this package. Citations: the retained excerpts carry no citation command at all, so no bibliography entry of the article is transcribed and no reference is redistributed. Reference adaptation: none was needed and none was made; every reference inside the delivered unit resolves inside it, so no unresolved reference or citation is produced. Printed slips kept literal: no printed word, symbol, hypothesis or number of the counted statement or of its proof was altered. Layout and class: the article's own class is \texttt{quantumarticle} with packages including braket, amsmath, amsfonts, amsthm, amssymb, url, graphicx, xcolor, hyperref, subcaption, appendix, comment, natbib; the wrapper is the lane's pinned \texttt{amsart} editorial wrapper at 10pt on A4, which restates the theorem-like environments the retained text needs and adds the article's own macro definitions that the retained text needs (\texttt{\textbackslash E}, \texttt{\textbackslash L}, \texttt{\textbackslash NN}, \texttt{\textbackslash pl}), with extra wrapper packages (bm, bbm, dsfont, xspace, cancel, mathtools). The delivered numbering is the unit's own. Assets: no retained span contains a figure environment, a graphics-inclusion command, a PDF-page inclusion, a table or a TikZ picture; no image file is copied into this package, the package declares no asset dependency and no image file is redistributed with it. Licence: version arXiv:2203.03745v5 of this article is distributed under the Creative Commons Attribution 4.0 International licence, as recorded in the arXiv licence metadata for this exact version and shown on the abstract page https://arxiv.org/abs/2203.03745v5. A full rights notice is delivered at licenses/arxiv-2203.03745-CC-BY-4.0.txt. Characters: every retained excerpt is pure ASCII, so the delivered engine, XeLaTeX with its default Latin Modern text font, reports no missing character.
\begin{lemma} \label{lem:zenolem}
Let $(\Phi_m)_{m=1}^k$ be a family of contractions for any $k \in \NN$. Let $(\L_m)_{m=1}^k$ be bounded Lindbladians such that $e^{-\L_m}$ is also contractive for each $m$. Let $\|\L\| := \max_m \{\| \L_m \|\}$. Assume for given $g \in \NN$, projection $\E$, and $\epsilon \in (0,1)$ that for any contiguous indices $1 \leq j_1, ..., j_n \leq k$, $\Phi_{j_1} ... \Phi_{j_n} \geq_{cp} (1 - \epsilon) \E$, and that $\Phi_j \E = \E \Phi_j = \E$ for each $j \in 1...k$. Then
\begin{equation*}
\begin{split}
& \bigg \| \sum_{m=0}^k \frac{(-1)^m}{k^m} \sum_{W \in WS(m+1,k)} \Big ( \Big ( \prod_{j=1}^m (\E \L_{W,j})
    - \prod_{j=1}^m (\Phi_{W(j)} \L_{W,j}) \Big ) \Phi_{W(m+1)} \Big ) \E \bigg \|
\\ & \leq \frac{g}{k} \Big ( 1 + \frac{1}{k} \Big ) \frac{1}{1-\epsilon} (\|\L\|^3 + 4 \|\L\|^2 + 2 \|\L\|) e^{\|\L\|}
\end{split}
\end{equation*}
where $\Phi_{W(j)} = \Phi_{W(j)[|W(j)|]} \circ ... \circ \Phi_{W(j)[1]}$, and each $\L_{W,j} \in (\L_m)$ is the Lindbladian appearing between the channel compositions $W(j)$ and $W(j+1)$.
\end{lemma}

\begin{proof}
For each $n \in 1...k$, let $\#(n, WS(m+1,k))$ denote the number of partitions in $WS(m+1,k)$ that contain at least one interval of length exactly $n$. To estimate $\#(n, WS(m+1,k))$, consider placing $m-1$ partition boundaries (creating $m$ partitions), then placing a final partition boundary that is exactly $n$ positions away from one of those already placed. To place the first $1...k$ boundaries has $(k+1 \text{ choose } m-1)$ choices when including placing before the 1st or after the $k$th as well as between any indices. To place the final boundary has 2 choices for each already placed boundary, although these terms are subsequently divided by a factor of 2 to account for exchange between the final and $n$-positions-away pre-existing boundary, plus another 2 choices for placement at either end of the interval, yielding a $m+1$ possibilities. The estimate $\#(n, WS(m+1,k)) \leq (k+1 \text{ choose } m-1) \times (m+1)$ is an overcount, since it may double-count if there was already a length-$n$ partition between the first $m-1$ boundaries, and since placing a boundary $n$ positions away from another may not create an $n$-length interval if another boundary is in between. To simplify, as long as $m > 1$,
\begin{equation} \label{eq:termnum}
    \#(n, WS(m+1,k)) \leq \binom{k+1}{m-1} (m+1) = \frac{(k+1) k! (m+1)}{(m-1)!(k - m + 2)!} \leq \frac{(k+1)k^{m-2} (m+1)}{(m-1)!} \pl.
\end{equation}
Futhermore, $\#(n, WS(m+1,k))$ overcounts the number of terms containing one partition of length $n$ and no smaller partitions.

Assuming that each $|W_j| \geq n$,
\begin{equation*}
\begin{split}
    & \Big ( \prod_{j=1}^m (\Phi_{W(j)} \L_{W,j}) \Big ) \Phi_{W(m+1)}
    \\ & = \prod_{j=1}^m (((1-\epsilon^{\lfloor n / g \rfloor}) \E + \epsilon^{\lfloor n / g \rfloor}  \Psi_{W(j)}) \L_{W,j}) ((1-\epsilon^{\lfloor n / g \rfloor}) \E + \epsilon^{\lfloor n / g \rfloor} \Psi_{W(m+1)} )
\end{split}
\end{equation*}
for some contractions $(\Psi_{W(j)})_{j = 1}^{m+1}$ such that $\E \Psi_{W(j)} = \Psi_{W(j)} \E = \E$. Hence
\begin{equation} \label{eq:singleterm}
\Big \| \Big ( \prod_{j=1}^m (\Phi_{W(j)} \L_{W,j}) \Big ) \Phi_{W(m+1)} \E - \prod_{j=1}^m (\E \L_{W,j}) \E \Big \|
    \leq \Big ( \prod_{j=1}^m \|\L_{W,j}\| \Big ) \big (1 - (1-\epsilon^{\lfloor n / g \rfloor})^{m} \big ) \pl,
\end{equation}
 assuming that the smallest partition in $W$ is of length at least $n$.
 %Via the inequality of arithmetic and geometric means,
 %\[ \prod_{j=1}^m \|\L_{W,j}\| \leq \Big ( \frac{1}{m} \sum_j \|\L_{W,j}\| \Big )^m \pl. \]
 Via Bernoulli's inequality, $(1 - (1-\epsilon^{\lfloor n / g \rfloor})^m \big ) \leq m \epsilon^{\lfloor n / g \rfloor}$ when $m > 0$.

To bound the full expression as in the Lemma,
\begin{equation*}
\begin{split}
& \bigg \| \sum_{m=0}^k \frac{(-1)^m}{k^m}
    \bigg (  \sum_{W \in WS(m+1,k)} \bigg (  \Big ( \prod_{j=1}^m (\Phi_{W(j)} \L_{W,j}) \Big ) \Phi_{W(m+1)} \E - \prod_{j=1}^m (\E \L_{W,j}) \E \big ) \bigg ) \bigg ) \bigg \|
\\ & \leq \sum_{m=0}^k \frac{(-1)^m}{k^m} \sum_{W \in WS(m+1,k)}
      \bigg \| \bigg (  \Big ( \prod_{j=1}^m (\Phi_{W(j)} \L_{W,j}) \Big ) \Phi_{W(m+1)} \E - \prod_{j=1}^m (\E \L_{W,j}) \E \big ) \bigg ) \bigg \|
\\ & \leq \sum_{m=0}^k \frac{(-1)^m}{k^m} \sum_{n=1}^k \sum_{W \in WS(m+1,k,n)}
      \bigg \| \bigg (  \Big ( \prod_{j=1}^m (\Phi_{W(j)} \L_{W,j}) \Big ) \Phi_{W(m+1)} \E - \prod_{j=1}^m (\E \L_{W,j}) \E \big ) \bigg ) \bigg \|
    \pl,
\end{split}
\end{equation*}
where $W \in WS(m+1,k,n)$ counts all partionings with minimum inveral size $n$. Combining with Equations \eqref{eq:termnum} and \eqref{eq:singleterm},
\begin{equation} \label{eq:zlemmamain}
\begin{split}
& \bigg \| \sum_{m=0}^k \frac{(-1)^m}{k^m} \bigg (  \sum_{W \in WS(m+1,k)} \bigg (  \Big ( \prod_{j=1}^m (\Phi_{W(j)} \L_{W,j}) \Big ) \Phi_{W(m+1)} \E - \prod_{j=1}^m (\E \L_{W,j}) \E \big ) \bigg ) \bigg ) \bigg \|
\\ & \leq \Big ( \frac{1}{k} + \frac{1}{k^2} \Big ) \Big ( \sum_{n=0}^k  \epsilon^{\lfloor n / g \rfloor} \Big ) \sum_{m=1}^k \|\L \|^m\frac{(m+1) m}{(m-1)!} \pl.
\end{split}
\end{equation}
Estimating the first sum,
\begin{equation} \label{eq:nsum}
    \sum_{n=0}^k  \epsilon^{\lfloor n / g \rfloor} \leq g \sum_{n=0}^{\lceil k/g \rceil } \epsilon^n \leq \frac{g}{1-\epsilon} \pl.
\end{equation}
For the second sum,
\begin{equation} \label{eq:msummain}
    \sum_{m=1}^k \|\L \|^m \frac{(m+1) m}{(m-1)!} = \sum_{m=1}^k \|\L \|^m \frac{m^2}{(m-1)!} + \sum_{m=1}^k \|\L \|^m \frac{m}{(m-1)!} \pl.
\end{equation}
Note that whenever a sum starting at $m=0$ contains a factor of $m$ in all terms, it equivalently starts at $m=1$. The first term in Equation \eqref{eq:msummain} decomposes as
\begin{equation} \label{eq:msumterm1}
\begin{split}
& \sum_{m=1}^k \|\L \|^m\frac{m^2}{(m-1)!} = \|\L\| \sum_{m=0}^\infty \frac{(m+1)^2}{m!} \|\L\|^m
    \\ & = \|\L\| \Big ( \sum_{m=1}^k \frac{m^2}{m!} \|\L\|^m + 2 \sum_{m=1}^k \frac{m}{m!} \| \L \|^m + \sum_{m=0}^k \frac{1}{m!} \| \L \|^m \Big ) \pl.
\end{split}
\end{equation}
The second term in Equation \eqref{eq:msummain} is equivalent to
\begin{equation} \label{eq:msumterm2}
\begin{split}
\sum_{m=1}^k \|\L \|^m \frac{m}{(m-1)!} & = \| \L \| \sum_{m=0}^k \frac{m+1}{m!} \| \L \|^m
         = \| \L \| \sum_{m=0}^k \frac{m}{m!} \| \L \|^m + \| \L \| \sum_{m=0}^k \frac{1}{m!} \| \L \|^m \pl,
\end{split}
\end{equation}
which also happens to equal the first term in parenthesis in the last line of Equation \eqref{eq:msumterm1}. Moreover,
\begin{equation} \label{eq:onem}
    \sum_{m=1}^k \frac{m}{m!} \|\L\|^m =  \sum_{m=1}^k \frac{1}{(m-1)!} \| \L\|^m
        \leq \|\L\| \sum_{m=0}^\infty \frac{1}{m!} \|\L\|^m = \|\L\| e^{\|\L\|} \pl.
\end{equation}
Therefore, the terms in Equation \eqref{eq:msumterm2} are equivalent to $(\| \L \|^2 + \| \L \|) e^{\| \L \|}$. Collecting terms from Equation \eqref{eq:msumterm1} yields $(\|\L\|^3 + 3 \| \L \|^2 + \|\L\|) e^{\|\L\|}$. Adding these expressions yields
\begin{equation} \label{eq:msumfinal}
    \sum_{m=1}^k \|\L \|^m\frac{m^2}{(m-1)!} \leq (\|\L\|^3 + 4 \|\L\|^2 + 2 \|\L\|) e^{\|\L\|} \pl.
\end{equation}
Equations \eqref{eq:nsum} and \eqref{eq:msumfinal} with \eqref{eq:zlemmamain} complete the Lemma.
\end{proof}

\end{document}
